% =============================================================================
% 6.6 GESTIÓN DE TRANSACCIONES Y CONTROL DE CONCURRENCIA
% =============================================================================
\section{Gestión de Transacciones y Control de Concurrencia}
\label{sec:concurrencia}

PostgreSQL emplea \textbf{MVCC} (\textit{Multi-Version Concurrency Control}): los lectores
nunca bloquean a los escritores y viceversa, ya que cada transacción ve una instantánea
consistente de los datos. EthosHub se apoya en este modelo y lo complementa con
estrategias específicas en la capa de persistencia.

\subsection{Niveles de aislamiento transaccional}
\label{subsec:aislamiento}

El backend opera bajo el nivel de aislamiento por defecto de PostgreSQL,
\comando{READ COMMITTED}, adecuado para la mayoría de las operaciones CRUD. La
tabla~\ref{tab:aislamiento} resume las garantías relevantes.

\vspace{0.2cm}
\noindent
\renewcommand{\arraystretch}{1.3}
\fontsize{9}{10.8}\selectfont\sffamily
\begin{tabularx}{\textwidth}{|>{\bfseries\color{colorPrimario}}l|X|}
\hline
\rowcolor{headerTabla}
\textbf{Aspecto} & \textbf{Comportamiento en EthosHub} \\
\hline
Nivel por defecto & \texttt{READ COMMITTED}: cada sentencia ve los datos confirmados al iniciarse. \\
\hline
\rowcolor{grisTecnico}
Idempotencia & El índice único parcial de \texttt{chat\_messages} evita duplicados ante reintentos concurrentes. \\
\hline
Operaciones atómicas & \texttt{UPSERT} (\texttt{INSERT ... ON CONFLICT}) en \texttt{add\_like}, \texttt{record\_login\_failure}, \texttt{upsert\_chat}. \\
\hline
\rowcolor{grisTecnico}
Pooling & El acceso pasa por el \textit{pooler} de Supabase, acotando conexiones simultáneas (RNF-CONC-01). \\
\hline
\end{tabularx}
\label{tab:aislamiento}

\subsection{Control transaccional desde la capa de persistencia}
\label{subsec:control_tx}

Desde el backend, la capa tipada con \textbf{jOOQ} y \textbf{Spring JDBC} delimita las
transacciones; las operaciones compuestas (p.~ej. crear un perfil y su configuración)
se ejecutan dentro de una única transacción gestionada por Spring
(\texttt{@Transactional}). Las funciones \texttt{SECURITY DEFINER} que realizan varias
escrituras (como \comando{delete\_profile\_cascade}) constituyen, por sí mismas, una
unidad transaccional atómica en el motor.

\begin{NotaTecnica}
El patrón \comando{INSERT ... ON CONFLICT DO NOTHING/UPDATE} se prefiere a la secuencia
``leer-decidir-escribir'' porque resuelve la condición de carrera \textbf{dentro de una
única sentencia atómica}, eliminando ventanas de tiempo en las que dos transacciones
concurrentes podrían insertar duplicados.
\end{NotaTecnica}
